\documentclass[11pt,a4paper]{article}
\usepackage[T1]{fontenc}
\usepackage[utf8]{inputenc}
\usepackage{lmodern}
\DeclareUnicodeCharacter{1E93}{\d{z}}
\usepackage{amsmath,amssymb}
\usepackage[margin=25mm]{geometry}
\usepackage[hidelinks]{hyperref}
\hypersetup{pdftitle={Quantum wall dynamics: boundary choice, reflected packets and impulse},pdfauthor={navstokgap research working notes}}
\newcommand{\tightlist}{\setlength{\itemsep}{0pt}\setlength{\parskip}{0pt}}
\setlength{\emergencystretch}{3em}
\setcounter{secnumdepth}{0}
\begin{document}
\section{Quantum wall dynamics: boundary choice, reflected packets and
impulse}\label{quantum-wall-dynamics-boundary-choice-reflected-packets-and-impulse}

\textbf{Result, 2026-10-27.} For the specified incident Gaussian and
free Dirichlet half-line dynamics below, the nonnegative boundary-force
measures converge weakly on the whole time axis: \[
 \frac{\hbar^2}{2m}\,|\partial_q\psi_{\hbar,D}(0,t)|^2\,dt
 \;\Longrightarrow\;2|p_0|\,\delta_{t_*},
 \qquad t_*=-\frac{mq_0}{p_0}>0.\tag{1}
\] This includes the impact layer excluded from a uniform off-impact
momentum-current estimate. In fact, on the \(\sqrt\hbar\) time scale the
force has an explicit Gaussian profile and the momentum current tends to
an error-function transition.

For the corresponding \textbf{Neumann} half-line packet, the derivative
in (1) vanishes identically. The complete boundary stress has another
term; it has the same limiting impulse and the same leading impact-layer
profile. The particle's classical bounce therefore does not determine
whether the quantum wave satisfies Dirichlet or Neumann boundary
conditions. General Robin conditions are discussed through their
scattering phases; no wavepacket limit is claimed here for arbitrary
Robin parameters.

For a particle falling back to a mirror, the time spent at different
heights also measures how many trajectories turn in a chosen interval.
One bounded average of recorded heights can give this count a negative
sign. In linear gravity the size of such a deficit would bound a trial
action even if the initial state changes with the action.
\hyperref[counting-turning-heights-from-neutron-positions]{Turning
heights and a finite clock}.

The boundary contribution to Ehrenfest's theorem is established
literature: see
\href{https://www.scielo.org.mx/scielo.php?pid=S1870-35422013000200001&script=sci_arttext}{De
Vincenzo (2013)}, equations (38)--(39) for a wall on the opposite
half-line (passage read), and
\href{https://arxiv.org/abs/2206.07531}{Albrecht et al.~(2023)}
(abstract) on boundary momentum and bouncing packets. A particularly
close precedent is
\href{https://arxiv.org/pdf/quant-ph/0408182}{Belloni--Doncheski--Robinett
(2004)}, equations (8)--(11), (16)--(19): they normalize the same
odd-image Gaussian on the opposite half-line and analyze its
near-collision position and force (relevant passages read). The
computation below gives a particular weak impulse measure and compares
the Neumann stress; it does not claim priority for the reflected
Gaussian or wall forces.

\subsection{Classical impact and quantum boundary
data}\label{classical-impact-and-quantum-boundary-data}

Take the half-line \(q\ge0\) and a free particle starting at \(q_0>0\)
with \(p_0<0\). For an ideal \textbf{elastic} wall, the prescribed
impact law is \(p(t_*+)=-p(t_* -)\) at \(t_*=-mq_0/p_0\). The motion is
\[q(t)=|q_0+p_0t/m|,\qquad
 p(t)=p_0\operatorname{sgn}(q_0+p_0t/m)\quad(t\ne t_*).\] Position is
continuous; momentum has a jump. A smooth-force equation with a finite
smooth force does not describe that jump. Nonpenetration alone does not
select the elastic law: a restitution coefficient or another contact
rule would be additional data.

For a smooth wall in several dimensions, the elastic specular rule
reverses the normal component and retains the tangential component:
\(v_+=v_--2(v_-\cdot n)n\). This requires a specified unit normal and a
well-defined collision. Corners, grazing contacts and an accumulation of
impacts require separate analysis; one should not infer their uniqueness
merely from the local smooth-force theorem.

A mathematically defined quantum impenetrable wall can be represented by
the Dirichlet Laplacian
\[H_{\hbar,D}=-\frac{\hbar^2}{2m}\frac{d^2}{dq^2},
\qquad D(H_{\hbar,D})=H^2(0,\infty)\cap H^1_0(0,\infty).\] It is
self-adjoint and gives a unitary group. Its Euclidean kernel is \[
 K_D(T;x,y)=k_{\hbar,T}(x-y)-k_{\hbar,T}(x+y),
 \quad
 k_{\hbar,T}(r)=\sqrt{\frac m{2\pi\hbar T}}
             e^{-mr^2/(2\hbar T)},\quad x,y>0.
\] This is the method of images: the Euclidean reference process is
Brownian motion killed at the boundary. The quantum real-time wall
reflects, but killed Euclidean paths should not themselves be read as
real-time bouncing trajectories. The reflected-packet estimate below
proves a quantitative coherent-packet position limit through the bounce.

A subtle domain issue is important: the symmetric differential operator
\(-i\hbar\,d/dq\) on the half-line with Dirichlet boundary has unequal
deficiency indices and no self-adjoint extension on \(L^2(0,\infty)\).
The Dirichlet Hamiltonian is nevertheless self-adjoint. A current or a
quadratic-form expectation involving \(-i\hbar d/dq\) can be studied,
but it should not be presented as the expectation of a self-adjoint
half-line momentum operator. See
\href{https://arxiv.org/abs/2103.01715}{Al-Hashimi and Wiese,
arXiv:2103.01715} (abstract) for the boundary momentum problem.

\subsection{1. The packet and its boundary
value}\label{the-packet-and-its-boundary-value}

Fix \(m,\omega,q_0>0\) and \(p_0<0\), with \(v=p_0/m\). On the full line
take the normalized coherent state \[
 g_{\hbar,0}(q)=\left(\frac{m\omega}{\pi\hbar}\right)^{1/4}
  \exp\!\left[-\frac{m\omega(q-q_0)^2}{2\hbar}
       +\frac{i p_0(q-q_0)}{\hbar}\right].
\] Let \(g_{\hbar,t}=e^{-itP^2/(2m\hbar)}g_{\hbar,0}\), \(a(t)=q_0+vt\),
\(A(t)=1+\omega^2t^2\), and \[
 \gamma_\hbar=\langle g_{\hbar,0},Rg_{\hbar,0}\rangle
 =\exp\!\left[-\frac{m\omega q_0^2+p_0^2/(m\omega)}{\hbar}\right],
 \qquad (Rg)(q)=g(-q).\tag{2}
\] The odd and even extensions restrict to the normalized half-line
states \[
 \psi_{\hbar,D}(q,t)
   =\frac{g_{\hbar,t}(q)-g_{\hbar,t}(-q)}{\sqrt{1-\gamma_\hbar}},
 \qquad
 \psi_{\hbar,N}(q,t)
   =\frac{g_{\hbar,t}(q)+g_{\hbar,t}(-q)}{\sqrt{1+\gamma_\hbar}},
 \quad q>0.\tag{3}
\] They solve the free Schrödinger equation with Dirichlet and Neumann
conditions, respectively. Their regularity and decay justify the
boundary calculations below.

The free Gaussian formula gives \[
 |g_{\hbar,t}(0)|^2
 =\sqrt{\frac{m\omega}{\pi\hbar A(t)}}
    \exp\!\left[-\frac{m\omega a(t)^2}{\hbar A(t)}\right],\qquad
 \hbar\,\frac{g'_{\hbar,t}(0)}{g_{\hbar,t}(0)}
 =D(t):=\frac{m\omega q_0+i p_0}{1+i\omega t}.\tag{4}
\] The simplification in \(D\) follows from \(a=q_0+p_0t/m\);
\(g_{\hbar,t}(0)\) is never zero. Define \(B=(m\omega q_0)^2+p_0^2\).
Since \(\psi'_{\hbar,D}(0,t)=2g'_{\hbar,t}(0)/\sqrt{1-\gamma_\hbar}\),
the boundary-force density is exactly \[
 F_{\hbar,D}(t)
 =\frac{2B}{m(1-\gamma_\hbar)A(t)}
    \sqrt{\frac{m\omega}{\pi\hbar A(t)}}
    \exp\!\left[-\frac{m\omega a(t)^2}{\hbar A(t)}\right].\tag{5}
\]

\subsubsection{Position through impact and current away from
impact}\label{position-through-impact-and-current-away-from-impact}

Write \(a(t)=q_0+p_0t/m\), and let \(g_{\hbar,t}\) be the freely evolved
Gaussian. Its position variance is
\[\sigma_q(t)^2=\frac{\hbar}{2m\omega}(1+\omega^2t^2).\tag{P21}\]
Because \(||x|-|a||\le|x-a|\) and reflection commutes with the free
evolution, \[
 \left|\langle Q\rangle_{\psi_\hbar(t)}-|a(t)|\right|
 \le\sqrt{\frac{2}{1-\gamma_\hbar}}\,\sigma_q(t).\tag{P22}
\] Indeed the half-line position expectation equals the full-line
expectation of \(|X|\) in the normalized full-line odd state
\(\phi_\hbar=(g_\hbar-Rg_\hbar)/\sqrt{2(1-\gamma_\hbar)}\); apply the
triangle inequality to the two Gaussian terms. Thus position converges
uniformly on every bounded time interval, \textbf{including the impact
time}, to the elastic classical folded trajectory.

Away from the impact, the corresponding probability-current momentum
also converges to the incoming or reflected value. More precisely, if
\(|a(t)|\ge\eta>0\), with \(\sigma_p=\sqrt{\hbar m\omega/2}\), \[
 \left|\int_0^\infty\overline{\psi_\hbar}
             (-i\hbar\partial_q)\psi_\hbar\,dq
       -p_0\operatorname{sgn}a(t)\right|
 \le\sqrt{\frac2{1-\gamma_\hbar}}
      \left(\sigma_p+\frac{2|p_0|\sigma_q(t)}{\eta}\right).\tag{P23}
\] For the proof, the half-line integral is the full-line expectation of
\(\operatorname{sgn}(X)P\) in the normalized full-line odd state
\(\phi_\hbar=(g_\hbar-Rg_\hbar)/\sqrt{2(1-\gamma_\hbar)}\). On the
Gaussian, \(\|(\operatorname{sgn}X\,P-\operatorname{sgn}a\,p_0)g_t\|\)
is at most \(\sigma_p+2|p_0|\mathbb P(aX<0)^{1/2}\); Chebyshev bounds
the last probability by \(\sigma_q(t)^2/a(t)^2\). The reflected Gaussian
obeys the same estimate. For each fixed \(T,\eta>0\) independent of
\(\hbar\), (P23) proves uniform convergence over \(|t|\le T\) with
\(|a(t)|\ge\eta\). It gives no uniform bound in an \(\hbar\)-dependent
impact layer \(|a(t)|\sim\sqrt\hbar\) or momentum value at \(t_*\). The
expression in (P23) is a real current form on this domain, not an
expectation of a self-adjoint half-line momentum.

Position converges uniformly through impact; current cannot converge
uniformly to a discontinuous momentum jump. The next calculation
resolves that layer in the weak measure topology.

\subsection{2. The pulse and the limiting
impulse}\label{the-pulse-and-the-limiting-impulse}

Put \(A_*=A(t_*)\) and
\[\Lambda=|v|\sqrt{\frac{m\omega}{A_*}}>0.\tag{6}\] Because
\(m\omega q_0=-p_0\omega t_*\), one has \(B=p_0^2A_*\). At every fixed
\(s\in\mathbb R\), (5) gives \[
 \sqrt\hbar\,F_{\hbar,D}(t_*+\sqrt\hbar\,s)
 \longrightarrow
 \frac{2p_0^2}{m}\sqrt{\frac{m\omega}{\pi A_*}}\,
          e^{-\Lambda^2s^2}.\tag{7}
\] The integral of the right side is \(2|p_0|\).

This pointwise calculation proves (1) with the following tail control.
For any fixed \(\delta>0\), \[
 c_\delta:=\inf_{|t-t_*|\ge\delta}\frac{m\omega a(t)^2}{A(t)}>0:
\] the expression is continuous, has only the one zero \(t_*\), and
tends to \(mv^2/\omega>0\) at both infinities. Equation (5) bounds the
total force outside this neighbourhood by
\(C_\delta\hbar^{-1/2}e^{-c_\delta/\hbar}\), since \(A(t)^{-3/2}\) is
integrable. Inside the neighbourhood, after \(t=t_*+\sqrt\hbar s\), the
profiles have a common integrable Gaussian majorant. Dominated
convergence proves convergence against every bounded continuous test
function and gives the total mass \(2|p_0|\). This is weak convergence
of measures, not total-variation convergence to an atom.

The real current-momentum form is \[J_\hbar(t)=\hbar\operatorname{Im}
 \int_0^\infty\overline{\psi_\hbar}\,\partial_q\psi_\hbar\,dq.\tag{8}\]
The differential operator \(-i\hbar\partial_q\) on a half-line is not a
self-adjoint momentum operator; (8) is a real current form. For a
sufficiently regular free half-line wave, \[
 \frac{dJ_\hbar}{dt}
 =\frac{\hbar^2}{2m}\left(
     |\psi_\hbar'(0,t)|^2
    -\operatorname{Re}[\overline{\psi_\hbar(0,t)}
                                  \psi_\hbar''(0,t)]\right).\tag{9}
\] To see this, differentiate (8), insert
\(\partial_t\psi=i\hbar\psi''/(2m)\), and integrate by parts. For
Dirichlet the second term vanishes, so \(dJ_{\hbar,D}/dt=F_{\hbar,D}\).
Using the off-impact current estimate above together with (7), \[
 J_{\hbar,D}(t_*+\sqrt\hbar\,s)
 \longrightarrow |p_0|\,\operatorname{erf}(\Lambda s).\tag{10}
\] The incoming value is \(p_0=-|p_0|\), the outgoing value is
\(|p_0|\), and their difference is the mass in (1). The transition is
continuous for each \(\hbar>0\) and becomes discontinuous only in the
limit.

\subsection{3. Neumann and Robin boundary
information}\label{neumann-and-robin-boundary-information}

For \(\psi_{\hbar,N}'(0,t)=0\), the expression in (1) is identically
zero. Equation (9), rather than the Dirichlet-only term, is the
appropriate boundary stress. Since \[
 \frac{g_{\hbar,t}''(0)}{g_{\hbar,t}(0)}
 =\frac{D(t)^2}{\hbar^2}
   -\frac{m\omega}{\hbar(1+i\omega t)},
\] it is explicitly \[
 F_{\hbar,N}(t)=\frac{2|g_{\hbar,t}(0)|^2}{m(1+\gamma_\hbar)}
 \left[-\operatorname{Re}D(t)^2+\frac{\hbar m\omega}{A(t)}\right].
 \tag{11}
\] This density need not be pointwise nonnegative. At the impact,
\(D(t_*)=ip_0\), and therefore the rescaled Neumann density has the same
limit (7). Moreover \[
 |F_{\hbar,N}(t)|
 \le\frac{2|g_{\hbar,t}(0)|^2}{m(1+\gamma_\hbar)}
      \left[\frac{B}{A(t)}+\frac{\hbar m\omega}{A(t)}\right],
\] so the same argument after (7) controls the total variation outside
every impact neighbourhood. Consequently the finite signed measures
\(F_{\hbar,N}(t)dt\) converge weakly to the same \(2|p_0|\delta_{t_*}\).
Their impact-layer current has the same limit as (10). The
even-reflection argument also gives the same folded classical position
\(|a(t)|\). Thus the limiting position and impulse do not recover the
Dirichlet/Neumann reflection phase.

All self-adjoint realizations of the free half-line Laplacian use
Dirichlet or a real Robin condition \(\psi'(0)=\kappa\psi(0)\),
including Neumann at \(\kappa=0\). For an incoming stationary wave
\(e^{-ikq}+r_\kappa(k)e^{ikq}\), \[
 r_\kappa(k)=\frac{ik+\kappa}{ik-\kappa},\qquad
 |r_\kappa(k)|=1,\quad
 r_0=1,\quad r_D=-1.\tag{12}
\] At fixed real \(\kappa\) and fixed nonzero physical momentum,
\(k=|p|/\hbar\to\infty\) and \(r_\kappa(k)\to1\). This is a stationary
scattering comparison, not a proof of the impulse limit for every Robin
packet. For \(\kappa<0\) there is additionally a bound state
\(\sqrt{-2\kappa}e^{\kappa q}\) of energy \(-\hbar^2\kappa^2/(2m)\),
showing another quantum distinction that the free incoming bounce does
not encode.

\subsection{Counting turning heights from neutron
positions}\label{counting-turning-heights-from-neutron-positions}

\emph{Historical experimental application before the 18 November scope
reset.}

\textbf{6.1 Sol Max, 2026-11-18.} A classical neutron that visits the
top of a gravitational arc must spend time at the heights below it. This
gives a test using one terminal position per neutron. It requires
neither a model for the distribution of launch energies nor a fit to an
interference pattern. Newton's 1713 Lemma X, Corollary 4, supplies the
mechanical calibration: measure the downward acceleration \(g\) from
initial displacements against a ruler and a clock. The observable below
is a modern application to that measured force.
\href{newton-insertion-action.md\#lemma-x-and-the-measurement-of-initial-force}{Lemma
X in 1687 and 1713} (historical passages).

\textbf{A count from heights.} Let \(z\) be height above the reflecting
surface. Choose a smooth nonnegative window \(\phi\) supported on an
interval \([b,b+d]\) strictly above it. Give \(\phi\) units of the
square root of length. At the recorded height assign the dimensionless
score

\[
 B_\phi(z)=\int_0^z\frac{\phi'(x)}{\sqrt{z-x}}\,dx .
\]

Set the score to zero below the support of the window. It is bounded,
with \(M=\|B_\phi\|_\infty\le
2\sqrt d\|\phi'\|_\infty\) and \(L=\|B_\phi'\|_\infty\le
2\sqrt d\|\phi''\|_\infty\). On a full excursion of apex height \(H\),
Newton's fall law gives \(dt=|dz|/\sqrt{2g(H-z)}\) on either leg. The
Abel integral \(\int_x^H dz/\sqrt{(z-x)(H-z)}=\pi\) therefore gives

\[
 \int_{\rm excursion} B_\phi(z(t))\,dt
       =\pi\sqrt{2/g}\,\phi(H)\ \ge 0 .
\]

Every return from below \(b\) obeys this identity, even if the wall
changes the energy between successive arcs. A partial excursion has
absolute integral at most \(\pi\sqrt{2/g}\|\phi'\|_1\), by using
\(|\phi'|\) in the calculation. At most two excursions meet the ends of
an interval of duration \(T\). Sampling a time uniformly in that
interval, independently of the neutron, gives a mean score at least
\(-C_{\rm cl}\), where \(C_{\rm cl}=2\pi\sqrt{2/g}\|\phi'\|_1/T\). The
same bound holds for the sum of two independent uniform waits on
\([0,T]\) by conditioning on the first wait. Loss, creation or an
unmeasured impulse above \(b\) requires its own allowance. The proof
holds for continuous trajectories at every subdivision of the arc.
\href{https://dlmf.nist.gov/5.12.E1}{Euler beta integral} (formula).

There is an exact \textbf{central-force} version. Fix the conserved
angular momentum \(J\), or record it for each release. Let
\(U_J(r)=V(r)+J^2/(2mr^2)\) be the independently calibrated effective
potential and put \(u(r)=[U_J(r)-U_J(r_0)]/(mg)\), using a fixed
reference acceleration \(g>0\). Assume \(u'>0\) from the lower gate to
the apex. Score \(u'(r)B_\phi(u(r))\) instead. The radial speed is
\(\sqrt{2g(H-u(r))}\) for \(H=u(r_{\max})\). The factor \(u'\) converts
\(dr\) to \(du\) and gives exactly the same count and endpoint
allowance. Use gates below the support of \(\phi\) and bound the
response error in this transformed score. Thus even a nonconstant
central force, with its centrifugal contribution, admits the classical
test.

\textbf{How a reduced action enters the same readout.} For the action
bound specialize to uniform \(g\) and the Dirichlet half-line
Hamiltonian \(H_a=-a^2\partial_z^2/(2m)+mgz\). The mirror boundary
condition is an additional physical input. Set
\(\ell=(a^2/(2m^2g))^{1/3}\), \(\sigma=\ell/2^{2/3}\) and let
\(-\lambda_n\) be the negative Airy zeros in order of increasing
\(\lambda_n\). The normalized density is
\(\mathrm{Ai}((z-h_n)/\ell)^2/(\ell D_n)\) with \(h_n=\ell\lambda_n\)
and \(D_n=\mathrm{Ai}'(-\lambda_n)^2\).
\href{https://arxiv.org/pdf/1007.1623}{Belloni and Robinett, §2}
(normalized eigenfunctions).

In that state the measured score is exactly

\[
 \langle B_\phi\rangle_n =
 \frac{1}{2\sqrt{\ell}D_n}
 \int_{\mathbb R}\phi(x)\,\frac1{\sigma}
        \mathrm{Ai}\!\left(\frac{x-h_n}{\sigma}\right)\,dx .
\]

To apply the known full-line Airy position marginal, take its two energy
arguments equal. In dimensionless variables it gives
\(\mathrm{Ai}(y)^2=(2^{2/3}\pi)^{-1}\int_0^\infty
\mathrm{Ai}(2^{2/3}y+t)\,dt/\sqrt t\). Applying the same Abel integral
once more gives the tail \(\int_x^\infty\rho_n(z)\,dz/\sqrt{z-x}
 =(2\sqrt\ell D_n)^{-1}\int_{(x-h_n)/\sigma}^\infty\mathrm{Ai}(t)\,dt\)
for \(x>0\). Integration against \(\phi'\) gives the displayed score.
Only the density on the positive half-line is used.
\href{https://arxiv.org/pdf/hep-th/0011137}{Curtright, Uematsu and
Zachos, Eqs. (51)--(53)} (Airy position marginal) and
\href{https://dlmf.nist.gov/9.5.E1}{NIST, Eq. 9.5.1} (Fourier
representation).

Define \(\widehat\phi(k)=\int e^{-ikx}\phi(x)\,dx\) and
\(A_3=(2\pi)^{-1}\int |k|^3|\widehat\phi(k)|\,dk<\infty\). The Airy
convolution differs from \(\phi(h_n)\) by at most \(\sigma^3 A_3/3\),
since its Fourier multiplier is a phase with argument of magnitude
\(\sigma^3|k|^3/3\). Also \(D_n\ge D_0=\mathrm{Ai}'(0)^2>0\) by the Airy
norm identity. Thus \textbf{every diagonal mixture}, including
populations chosen anew for each \(a\), has mean score at least
\(-\ell^{5/2}A_3/(24D_0)\). \href{https://dlmf.nist.gov/9.11.E5}{NIST,
Eq. 9.11.5} (Airy norm identity).

\textbf{A finite clock for coherent preparations.} Bound the prepared
vertical energy by \(E_*=mgH_*\), and write \(v_*=\sqrt{2E_*/m}\). The
two independent uniform waits specified above permit arbitrary
coherences in that band. Adjacent angular frequencies \(E_n/a\) there
are separated by at least \(\gamma=\pi g/v_*\). Indeed the Dirichlet
interval inequality between consecutive Airy zeros gives
\(\lambda_{n+1}-\lambda_n\ge\pi/\sqrt{\lambda_{n+1}}\), since the local
squared wave number on the interval is at most \(\lambda_{n+1}\).
Multiplying by \(mg\ell/a\) gives the stated spacing.

A frequency difference \(\omega\) contributes at most \(4/(\omega T)^2\)
after the two waits. For any density matrix with diagonal \(p_n\),
\(|\rho_{n,n+j}|\le\sqrt{p_n p_{n+j}}\) and their sum over \(n\) is at
most one for each \(j\). Every score matrix element is bounded by \(M\).
Summing both signs of \(j\) bounds the entire off-diagonal correction by
\(C_{\rm q}=4\pi^2 M/(3\gamma^2 T^2)=4M v_*^2/(3g^2T^2)\). This uses the
actual exposure distribution for every trial action.

\textbf{A finite ruler.} Suppose the recorded height is \(z+\xi\), where
the response error has the same law at each height, independently of the
launch, and \(|\xi|\le e_z<b\). Averaging this readout replaces
\(\phi(x)\) by \(\bar\phi(x)=\mathbb E\,\phi(x+\xi)\) exactly: the
integral defining \(B\) commutes with translations on the real line when
the window is extended by zero. The effective window remains nonnegative
and supported above the wall. Its variation, score norm and \(A_3\) can
only decrease. Thus blur of this kind preserves the classical sign and
the quantum bound. A measured position-dependent remainder instead costs
at most \(L\) times its mean absolute error.

Let \(s\) be a lower confidence bound on minus the observed mean, after
subtracting readout and preparation errors. A fraction \(f\) of
erroneous or missing scores costs at most \(2Mf\) for the complete
launch ensemble. A deviation \(\epsilon_t\) in total variation from the
specified clock law costs at most \(2M\epsilon_t\), with the bound valid
conditional on the launch. An energy leakage probability \(\eta<1\)
costs at most \(2M\sqrt\eta\) by projection and the trace-distance
bound. If \(s>C_{\rm cl}\), the continuous classical trajectories fail.
If also \(s>C_{\rm q}\), the same fixed-ruler and fixed-clock
observation gives

\[
 a\ \ge\ m\sqrt{2g}
   \left(\frac{24D_0(s-C_{\rm q})}{A_3}\right)^{3/5}
\]

within the stated Dirichlet dynamics. Under \(a\mapsto t a\) the band
projector and allowed preparation change, while \(m,g,H_*,T,b,d\), the
height response and the clock law remain physical inputs. The diagonal
error scales as \(t^{5/3}\). Neither bound contains a mesh.
\href{https://arxiv.org/abs/quant-ph/9712042}{Fuchs and van de Graaf,
Theorem 1} (trace distance).

The sign can be attained at finite resolution. For instance, take
\(b=1\,\mu{\rm m}\), \(d=2\,\mu{\rm m}\) and
\(\phi(x)=\sqrt d[4y(1-y)]^6\) for \(y=(x-b)/d\) in \([0,1]\), zero
elsewhere. Fix a height response supported within
\(e_z=0.5\,\mu{\rm m}\) and choose a reference length
\(\ell_0=6\,\mu{\rm m}\) by the ruler. At \(a_0=m\sqrt{2g\ell_0^3}\),
even the enlarged window from \(0.5\) to \(3.5\,\mu{\rm m}\) maps inside
the first negative lobe \((-\lambda_2,-\lambda_1)\) of the ground
state's Airy kernel. Let \(c_A\) be the minimum of
\(-\mathrm{Ai}((x-h_1)/\sigma)\) on that enlarged interval at
\(\ell=\ell_0\). The mean score is at most \(-s_0\) with
\(s_0=c_A\int\phi(x)\,dx/(2\sqrt{\ell_0}D_1\sigma)>0\) for every such
height response. Take \(H_*=100\,\mu{\rm m}\) and choose \(T\) at least
the larger of \(8\pi\sqrt{2/g}\|\phi'\|_1/s_0\) and
\((4v_*/g)\sqrt{M/s_0}\). Both clock allowances are then at most
\(s_0/4\). For a beam of independently selected horizontal speed
\(v_\parallel\), a detector plane at distance \(v_\parallel(t_1+t_2)\)
implements the two waits on fresh launches; changes to the preparation
or the time law across planes must be bounded. The reference prediction
uses a ground state. The test of a measured deficit allows arbitrary
populations and coherences within the calibrated energy band.
\href{https://dlmf.nist.gov/9.9\#T1}{NIST, Table 9.9.1} (first two
zeros).

\textbf{What existing neutron records establish.} A simple check uses
two adjacent equal-height bins. For a complete ballistic arc the lower
occupancy is at least \(\sqrt2-1\) times the upper occupancy.
Integrating \(1/\sqrt{H-z}\) in the two bins gives the minimum when the
apex is at the upper edge of the upper bin. In Nesvizhevsky et al.'s
direct track profile the bins centred at \(2\) and \(4\,\mu{\rm m}\)
contain roughly 324 and 346 events, well inside that inequality. The
apparent ratio 47/161 at \(-2\) and \(0\,\mu{\rm m}\) involves a bin
below the fitted mirror height. The authors excluded three edge points
in their fit. Ichikawa et al.'s later magnified profile includes
propagation from the guide through a cylindrical reflector. They report
\(p=0.715\) for their quantum fit and \(p=0.062\) for their classical
guide-loss model.
\href{https://arxiv.org/pdf/hep-ph/0502081}{Nesvizhevsky et al., §4 and
Fig. 15} and \href{https://arxiv.org/pdf/1304.1660}{Ichikawa et al.,
Fig. 3 and fit} (apparatus and experimental results).

There is an exact classical ambiguity at a single detection plane. For a
fixed flight time and any terminal height law \(f(dz)\) within the
energy ceiling, put all the neutrons at vertical rest with that height
law at the readout plane and transport the measure backwards by the
elastic gravitational flow to define the source. It gives the specified
direct height law, including any calibrated position blur. The
construction can be made conditional on horizontal flight speed.
Independent randomized readout after launch tests that explanation. The
cited experiments do not report the clocked height statistic.
\href{../docs/NeutronGravity_HeightMeasurements.md}{Source and graphical
readings}.

\subsection{4. Consequence for STATE}\label{consequence-for-state}

For a specified coherent packet, the Dirichlet force is a positive
boundary measure converging to Newton's impulsive momentum jump; the
full Neumann boundary stress has the same limit. A classical bounce does
not identify the boundary condition or reflection phase. This constructs
neither a general boundary-groupoid theorem nor a
positive-action-necessity result. For gravitational returns the
outstanding observation is a negative turning-height score with the
independent clock and complete counts.
\end{document}
